% Options for packages loaded elsewhere
\PassOptionsToPackage{unicode}{hyperref}
\PassOptionsToPackage{hyphens}{url}
\PassOptionsToPackage{dvipsnames,svgnames,x11names}{xcolor}
\documentclass[
  11pt,
]{article}
\usepackage{xcolor}
\usepackage[margin=27mm]{geometry}
\usepackage{amsmath,amssymb}
\setcounter{secnumdepth}{-\maxdimen} % remove section numbering
\usepackage{iftex}
\ifPDFTeX
  \usepackage[T1]{fontenc}
  \usepackage[utf8]{inputenc}
  \usepackage{textcomp} % provide euro and other symbols
\else % if luatex or xetex
  \usepackage{unicode-math} % this also loads fontspec
  \defaultfontfeatures{Scale=MatchLowercase}
  \defaultfontfeatures[\rmfamily]{Ligatures=TeX,Scale=1}
\fi
\usepackage{lmodern}
\ifPDFTeX\else
  % xetex/luatex font selection
  \setmainfont[]{Cambria}
  \setmonofont[]{Consolas}
\fi
% Use upquote if available, for straight quotes in verbatim environments
\IfFileExists{upquote.sty}{\usepackage{upquote}}{}
\IfFileExists{microtype.sty}{% use microtype if available
  \usepackage[]{microtype}
  \UseMicrotypeSet[protrusion]{basicmath} % disable protrusion for tt fonts
}{}
\makeatletter
\@ifundefined{KOMAClassName}{% if non-KOMA class
  \IfFileExists{parskip.sty}{%
    \usepackage{parskip}
  }{% else
    \setlength{\parindent}{0pt}
    \setlength{\parskip}{6pt plus 2pt minus 1pt}}
}{% if KOMA class
  \KOMAoptions{parskip=half}}
\makeatother
\usepackage{longtable,booktabs,array}
\newcounter{none} % for unnumbered tables
\usepackage{calc} % for calculating minipage widths
% Correct order of tables after \paragraph or \subparagraph
\usepackage{etoolbox}
\makeatletter
\patchcmd\longtable{\par}{\if@noskipsec\mbox{}\fi\par}{}{}
\makeatother
% Allow footnotes in longtable head/foot
\IfFileExists{footnotehyper.sty}{\usepackage{footnotehyper}}{\usepackage{footnote}}
\makesavenoteenv{longtable}
\setlength{\emergencystretch}{3em} % prevent overfull lines
\providecommand{\tightlist}{%
  \setlength{\itemsep}{0pt}\setlength{\parskip}{0pt}}
\usepackage{bookmark}
\IfFileExists{xurl.sty}{\usepackage{xurl}}{} % add URL line breaks if available
\urlstyle{same}
\hypersetup{
  pdftitle={WP-14 Literature Review --- Residue Stress Sweep (Run 03)},
  pdfauthor={The Privacy-is-Value Research Programme},
  colorlinks=true,
  linkcolor={black},
  filecolor={Maroon},
  citecolor={Blue},
  urlcolor={blue},
  pdfcreator={LaTeX via pandoc}}

\title{WP-14 Literature Review --- Residue Stress Sweep (Run 03)}
\author{The Privacy-is-Value Research Programme}
\date{2026-08-17}

\begin{document}
\maketitle

\section{Lit-review runtime · run 03 manifest (residue stress
sweep)}\label{lit-review-runtime-run-03-manifest-residue-stress-sweep}

\textbf{Purpose.} Run 02 (wf\_ba400906-67b) returned 0 VALIDATED / 5
MIRAGE and left five residual novel cores. Those residues are the claims
WP-14 now rests on, so this run stresses the residues themselves against
a widened, time-advanced corpus (2023-2026 window, searched 2026-08-17).

\subsection{Seat configuration}\label{seat-configuration}

{\def\LTcaptype{none} % do not increment counter
\begin{longtable}[]{@{}
  >{\raggedright\arraybackslash}p{(\linewidth - 6\tabcolsep) * \real{0.2500}}
  >{\raggedright\arraybackslash}p{(\linewidth - 6\tabcolsep) * \real{0.2500}}
  >{\raggedright\arraybackslash}p{(\linewidth - 6\tabcolsep) * \real{0.2500}}
  >{\raggedright\arraybackslash}p{(\linewidth - 6\tabcolsep) * \real{0.2500}}@{}}
\toprule\noalign{}
\begin{minipage}[b]{\linewidth}\raggedright
Seat
\end{minipage} & \begin{minipage}[b]{\linewidth}\raggedright
Count
\end{minipage} & \begin{minipage}[b]{\linewidth}\raggedright
Role
\end{minipage} & \begin{minipage}[b]{\linewidth}\raggedright
Sees
\end{minipage} \\
\midrule\noalign{}
\endhead
\bottomrule\noalign{}
\endlastfoot
sweep:* & 5 & multi-modal prior-art search (WEIS/econ venues; formal
preprints; law+standards; data-as-labour; shelf-life/erosion) & its
modality brief + the PVM position + the five residues + WebSearch \\
refute:* & 5 & per-residue DELEGATOR - hunts covering art, instructed to
default to coverage & the residue + sweep digest + WebSearch; never a
prover argument (D3) \\
judge & 1 & rules SURVIVES / NARROWED / COVERED per residue; assembles
bibliography additions & residues + refuter outputs + sweep items \\
\end{longtable}
}

\textbf{D4b honest limit (unchanged from run 02).} All seats run on the
same model weights; the separation is context isolation, not weight
diversity. The operating rule stands: an all-SURVIVES result would be
treated as a failed enforcement of the adversarial discipline. This run
returned 1 COVERED / 4 NARROWED / 0 SURVIVES across 5 residues.

\subsection{Run statistics}\label{run-statistics}

\begin{itemize}
\tightlist
\item
  Sweep items surfaced: 76 across 5 modalities
\item
  Refuter covering candidates: 48
\item
  Bibliography additions proposed: 74
\item
  Verdict spread: \textbf{1 COVERED / 4 NARROWED / 0 SURVIVES across 5
  residues}
\end{itemize}

\subsection{Judge headline}\label{judge-headline}

1 COVERED, 4 NARROWED, 0 SURVIVES: LR-01's Phi/R assessment-unit move is
anticipated outright (Nissim 2018, Abowd-Schmutte 2019, Census DSEP,
NIST SP 800-226, Epsilon Registry) and should be retired into
LR-02/LR-05; every other residue loses its per-leg firsts to named art
and survives only as a joint composition --- the priced
adversary-relative drifting-capital quadruple (LR-02), the
WTA-WTP/IEEE-7012 single-share-function identification (LR-03, strongest
residue: no game-theoretic treatment of MyTerms exists anywhere), the
Bajari-routed self-defeat falsifier landing on venue rent (LR-04,
hardest-narrowed: Mohades-Savona 2026 already states the rent-not-wage
conclusion by another route), and the Mosca-form reconstruction-clock
substitution carried into subject-owned valuation (LR-05).

\subsection{Outputs (this directory)}\label{outputs-this-directory}

\begin{itemize}
\tightlist
\item
  \texttt{residue\_verdicts\_03.md} - per-residue verdict, restated
  residual, basis, and the refuter's covering art.
\item
  \texttt{bibliography\_additions\_03.md} - new resolvable citations
  proposed for the WP-14 bibliography (A4 verifies before tier-A use).
\item
  \texttt{\_run\_data\_03.json} - full structured run output (sweeps,
  refutations, judge).
\end{itemize}

\newpage

\section{Residue verdicts: run-02 residues stress-tested (run
03)}\label{residue-verdicts-run-02-residues-stress-tested-run-03}

Verdict spread: \textbf{1 COVERED / 4 NARROWED / 0 SURVIVES across 5
residues}. Vocabulary: SURVIVES = the residue stands as stated; NARROWED
= partial coverage forces a smaller residue (restated below); COVERED =
the residue as stated is anticipated. Confidence labels follow run-02
practice; percentages are runtime-internal and do not enter tier-A
artefacts (GR-2).

\subsubsection{CTR-LR-01 --- COVERED (high,
82\%)}\label{ctr-lr-01-covered-high-82}

\textbf{Residual restatement (what now survives).} As stated --- a
governance-facing, independently verifiable mechanism property Phi, held
two-level-distinct from the per-subject bound R it gates, as the SINGLE
UNIT of a privacy assessment/valuation instrument --- the move is
anticipated: Nissim et al.~2018 conduct a compliance assessment over the
mechanism guarantee while holding a per-subject adversary-inference
bound as the gated requirement; Abowd-Schmutte 2019 make epsilon the
single unit of a social-choice valuation instrument; Census DSEP
practice governs a per-subject RECONSTRUCTION threat by deliberating and
setting epsilon; NIST SP 800-226 and the Epsilon Registry
(Dwork-Kohli-Mulligan 2019, OpenDP) institutionalize the mechanism
property as the assessed, comparable unit; the DP-auditing literature
(arXiv:2210.08643 line) supplies the independent-verifiability
qualifier; Kaissis 2023 / Biswas 2024-2025 supply the Phi-gates-R formal
bridge. What survives is not an independent claim: only the
instantiation where R is WP-14's time-drifting adversary-relative R(t)
and the instrument is a subject-side capital VALUATION --- that sliver
folds into CTR-LR-02/CTR-LR-05 and should be merged there, with
CTR-LR-01 retired as a separate novelty claim and rewritten as
positioning/related-work.

\textbf{Basis.} Five strong candidates from independent institutional
strands (legal formalization, economics, deployed census governance,
standards, registry) jointly span the claimed composition, including
reconstruction as the gated per-subject threat (Census 2020 DAS); the
refuter itself recommended COVERED and its absence report concedes
conjunct pairs co-occur in single works. The only uncovered triple is
exactly the LR-02/LR-05 subject-matter, not the Phi/R assessment-unit
move.

\textbf{Covering art found by the refuter seat.}

\begin{itemize}
\tightlist
\item
  {[}strong{]} Kobbi Nissim, Aaron Bembenek, Alexandra Wood, Mark Bun,
  Marco Gaboardi, Urs Gasser, David R. O'Brien, Salil Vadhan, Thomas
  Steinke --- Bridging the Gap between Computer Science and Legal
  Approaches to Privacy --- Harvard Journal of Law \& Technology 31(2)
  (2018) --- A governance/compliance-assessment methodology whose two
  steps are exactly the two-level separation: (1) formalize the legal
  requirement as a per-subject adversary-inference bound (the R-level
  object), (2) prove the mechanism-level guarantee (DP, the Phi-level
  object) implies it. The assessment verdict (FERPA compliance) is
  conducted over Phi alone while the per-subject bound is held distinct
  as the gated requirement --- a single work doing the two-level Phi/R
  separation AS the unit of a governance instrument.
  \url{https://www.semanticscholar.org/paper/Bridging-the-Gap-between-Computer-Science-and-Legal-Nissim-Bembenek/22067befe017dc5760be667bf0ba9cf48b2ab650}
\item
  {[}strong{]} John M. Abowd, Ian M. Schmutte --- An Economic Analysis
  of Privacy Protection and Statistical Accuracy as Social Choices ---
  American Economic Review 109(1); arXiv:1808.06303 (2019) --- A
  valuation instrument (marginal-cost-equals-marginal-benefit social
  choice) whose single choice unit is epsilon, the mechanism property,
  with epsilon's inferential-disclosure semantics (bound on adversary
  posterior odds about an individual) as the per-subject quantity it
  gates. Covers the `Phi as the single unit of a valuation instrument'
  branch directly, in an economics venue with governance intent.
  \url{https://arxiv.org/abs/1808.06303}
\item
  {[}strong{]} NIST --- Guidelines for Evaluating Differential Privacy
  Guarantees --- NIST SP 800-226 (final) (2025) --- An institutional,
  governance-facing assessment instrument whose unit of evaluation is
  the mechanism guarantee itself (including how to interpret what it
  bounds for individuals and how implementations can be independently
  checked) --- the closest standards-body embodiment of `assess Phi, not
  per-subject outcomes'.
  \url{https://nvlpubs.nist.gov/nistpubs/SpecialPublications/NIST.SP.800-226.pdf}
\item
  {[}strong{]} U.S. Census Bureau (Abowd et al.) --- The modernization
  of statistical disclosure limitation at the U.S. Census Bureau (2020
  DAS: DSEP committee sets the privacy-loss budget in response to the
  reconstruction attack) (2020) --- A deployed governance process in
  which a per-subject reconstruction threat (the 2010 census
  reconstruction attack --- the R-level event) is governed by
  deliberating and setting the mechanism property epsilon (the Phi-level
  unit) via a policy committee (DSEP/DRB). The assessment/decision unit
  is Phi; reconstruction risk is what it gates. Institutional practice
  of the exact composition, minus the word `valuation'.
  \url{https://www.census.gov/content/dam/Census/library/working-papers/2020/adrm/The\%20modernization\%20of\%20statistical\%20disclosure\%20limitation\%20at\%20the\%20U.S.\%20Census\%20Bureau.pdf}
\item
  {[}strong{]} Cynthia Dwork, Nitin Kohli, Deirdre Mulligan ---
  Differential Privacy in Practice: Expose your Epsilons! --- Journal of
  Privacy and Confidentiality 9(2); plus the live OpenDP Differential
  Privacy Deployments Registry and NIST IR registry-governance report
  (2025) (2019) --- Positions the mechanism property epsilon as the
  public, comparable, oversight-enabling governance unit across
  deployments (the Epsilon Registry), explicitly for accountability
  pressure --- Phi as the single governance-facing assessment unit, now
  institutionally realized at registry.opendp.org.
  \url{https://journalprivacyconfidentiality.org/index.php/jpc/article/view/689}
\item
  {[}partial{]} Arpita Ghosh, Aaron Roth --- Selling Privacy at Auction
  --- EC 2011 / Games and Economic Behavior; arXiv:1011.1375 (2011) ---
  A valuation/market instrument in which epsilon (the mechanism
  property) is literally the priced and procured unit, chosen from
  individuals' privacy valuations --- covering `Phi as the single unit
  of a valuation instrument', though the gated per-subject quantity is
  generic privacy loss rather than a named reconstruction bound.
  \url{https://arxiv.org/abs/1011.1375}
\item
  {[}partial{]} Justin Hsu, Marco Gaboardi, Andreas Haeberlen, Sanjeev
  Khanna, Arjun Narayan, Benjamin C. Pierce, Aaron Roth --- Differential
  Privacy: An Economic Method for Choosing Epsilon --- CSF 2014;
  arXiv:1402.3329 (2014) --- An economic instrument that converts
  per-participant expected harm (an adversary-relative, per-subject
  quantity) into the choice of the mechanism parameter epsilon --- the
  Phi/R direction of travel embedded in a valuation procedure, though it
  collapses the two levels into a calibration rather than holding Phi as
  an audited standing unit. \url{https://arxiv.org/abs/1402.3329}
\item
  {[}partial{]} Georgios Kaissis, Jamie Hayes, Alexander Ziller, Daniel
  Rueckert --- Bounding data reconstruction attacks with the hypothesis
  testing interpretation of differential privacy --- arXiv:2307.03928
  (with Riess et al.~arXiv:2402.12861 and Balle-Cherubin-Hayes IEEE S\&P
  2022 as siblings) (2023) --- Supplies the formal two-level structure
  the residue relies on: a mechanism-level DP property that provably
  gates a per-subject/per-record reconstruction bound. Not a governance
  or valuation instrument, but it removes any novelty in the Phi-gates-R
  separation itself. \url{https://arxiv.org/abs/2307.03928}
\item
  {[}partial{]} Matthew Jagielski et al.~lineage / (authors per arXiv)
  --- A General Framework for Auditing Differentially Private Machine
  Learning --- NeurIPS 2022; arXiv:2210.08643 (2022) --- The DP-auditing
  literature makes the mechanism property Phi empirically and
  independently verifiable --- covering the `independently verifiable'
  qualifier of the residue's Phi. \url{https://arxiv.org/pdf/2210.08643}
\item
  {[}partial{]} Ruwimal Y. Pathiraja, Jerome P. Reiter --- Setting the
  Privacy Budget in Differential Privacy by Bounding Adversaries' Odds
  of Learning Sensitive Information --- arXiv:2607.04004 (2026) ---
  Practitioner-facing procedure translating an adversary-posterior bound
  (R-level) into the epsilon setting (Phi-level) --- a governance-usable
  two-step bridge, though it collapses the levels into a calibration
  rather than keeping Phi as the standing audited unit.
  \url{https://arxiv.org/abs/2607.04004}
\end{itemize}

\textbf{Absence report (D2).} not\_found\_in\_scope(a single published
instrument that simultaneously (i) names a time-drifting,
adversary-relative per-subject RECONSTRUCTION bound R(t) as the gated
quantity, (ii) holds an independently auditable mechanism property Phi
two-level-distinct from it as the standing assessed unit, and (iii) is
framed as a privacy VALUATION/pricing instrument for data-as-capital.
Searched: economics-of-epsilon literature (Abowd-Schmutte AER 2019, Hsu
et al.~CSF 2014, Ghosh-Roth EC 2011, federated data-trading pricing),
epsilon-governance/registry literature (Dwork-Kohli-Mulligan 2019,
OpenDP Deployments Registry, NIST IR 2025), standards/PIA instruments
(NIST SP 800-226, CSRC PIA definitions, PIA tool guides),
legal-compliance formalizations (Nissim et al.~Harvard JOLT 2018, Is
Privacy Privacy? 2018), Census 2020 DAS governance and
reconstruction-attack record, DP auditing frameworks (arXiv:2210.08643,
2506.16666, 2509.08704, 2602.17454), and the sweep-digest
reconstruction-bound corpus (Balle 2022, Kaissis 2023, Riess 2024,
Biswas 2024/2025, Saeidian 2024, Pathiraja-Reiter 2026). This absence is
reported per D2 rules and is NOT evidence of novelty --- each conjunct
is individually covered and two conjuncts co-occur in single works
(Nissim 2018: (ii)+(a PIA-analog); Abowd-Schmutte/Census:
(iii)-adjacent+(reconstruction-motivated Phi governance)).)

\subsubsection{CTR-LR-02 --- NARROWED (moderate-high,
75\%)}\label{ctr-lr-02-narrowed-moderate-high-75}

\textbf{Residual restatement (what now survives).} Only the four-way
JOINT composition survives: a PRICED R(t) that is simultaneously (i)
adversary-relative in the reconstruction-robustness sense --- the object
itself belongs to Balle-Cherubin-Hayes 2022 / Kaissis 2023 / Riess 2024
and must be cited as source; (ii) drifting via adversary
background-information accumulation, explicitly distinguished from
DP-composition depletion (Zhang-Zhu 2021), valuation/reward drift (Xu
2017, CHRONOS 2026, Bo mean-field 2025), and aging-protects direction
(Age-Dependent DP 2022); (iii) held as structurally inalienable
subject-owned durable capital, answering Cofone's Calabresi-Melamed
objections to propertization; and (iv) with the observer-to-subject gap
decomposed as venue-exclusion/market-position rent shown NON-REDUCIBLE
to Acemoglu et al.'s data-externality account and distinct in mechanism
from Bergemann-Bonatti 2024's platform-position rent. No per-leg or
per-pair first-ness (drift+pricing, adversary-relative budget, drifting
durable capital, inalienable compensation) may be claimed anywhere in
the paper.

\textbf{Basis.} Five independent search scopes (\textasciitilde30
distinct queries) each found single-leg and pair coverage but no
joint-quadruple hit; refuter fetched the closest pair-threat
(arXiv:2101.04357) directly. Coverage of the legs is CONFIRMED by
resolvable art (Zhang-Zhu, Balle, Farboodi-Veldkamp, Cofone, Acemoglu,
Bergemann-Bonatti, Pathiraja-Reiter); absence of the joint composition
is scope-limited, and two formal rival decompositions of the same gap
(externality; platform rent) impose a real non-reducibility burden that
keeps this below SURVIVES.

\textbf{Covering art found by the refuter seat.}

\begin{itemize}
\tightlist
\item
  {[}partial{]} Tao Zhang, Quanyan Zhu, ``On the Differential Private
  Data Market: Endogenous Evolution, Dynamic Pricing, and Incentive
  Compatibility'', arXiv:2101.04357 (2021) (2021) --- Prices privacy
  loss that accumulates across periods via DP composition, with
  dynamically set per-period epsilon and compensation under IC/IR
  constraints. Defeats the claim's contrast `time-drifting vs static
  budget' as a standalone novelty: a priced, time-evolving DP loss
  already exists. Still adversary-independent epsilon, transactional
  (alienable), no rent decomposition --- quadruple untouched but the
  drift+pricing pair is anticipated.
  \url{https://arxiv.org/abs/2101.04357}
\item
  {[}partial{]} Vincent Conitzer, Curtis R. Taylor, Liad Wagman, ``Hide
  and Seek: Costly Consumer Privacy in a Market with Repeat Purchases'',
  Marketing Science 31(2):277-292 (2012) (2012) --- Mainstream-economics
  precedent where the observer's knowledge of the subject accumulates
  over time through repeat interactions and the subject pays a price
  (costly anonymity) to reset it --- a priced,
  observer-knowledge-drifting privacy model. Covers the intuition behind
  adversary-side accumulation driving privacy value over time, without a
  reconstruction bound, capital holding, or rent decomposition.
  \url{https://papers.ssrn.com/sol3/papers.cfm?abstract_id=2160895}
\item
  {[}partial{]} Lei Xu et al., ``Dynamic Privacy Pricing: A Multi-Armed
  Bandit Approach With Time-Variant Rewards'', IEEE Trans. Information
  Forensics and Security (2017) (2017) --- Data-market pricing of
  privacy where valuations are explicitly time-variant and prices adapt
  online --- prior art for `dynamic privacy pricing' as a phrase and
  mechanism; drift is on valuations/rewards, not on an
  adversary-relative reconstruction bound.
  \url{https://www.researchgate.net/publication/316236814_Dynamic_Privacy_Pricing_A_Multi-Armed_Bandit_Approach_With_Time-Variant_Rewards}
\item
  {[}partial{]} Borja Balle, Giovanni Cherubin, Jamie Hayes,
  ``Reconstructing Training Data with Informed Adversaries'', IEEE S\&P
  2022, arXiv:2201.04845 (2022) --- Origin of the
  adversary-parameterized reconstruction-robustness (ReRo) bound:
  adversary-relative reconstruction bounds as a formal object are known
  art (extended by Kaissis et al.~arXiv:2307.03928 and
  arXiv:2402.12861). Unpriced, static, no capital/rent --- covers the
  `adversary-relative reconstruction bound' leg alone.
  \url{https://arxiv.org/abs/2201.04845}
\item
  {[}partial{]} Daron Acemoglu, Ali Makhdoumi, Azarakhsh Malekian,
  Asuman Ozdaglar, ``Too Much Data: Prices and Inefficiencies in Data
  Markets'', AEJ: Microeconomics 14(4) (2022) (2022) --- Formal
  competing decomposition of the depressed subject-side price via the
  data/correlation externality --- an accumulating,
  other-subjects-relative information channel. The market-position-rent
  (venue-exclusion) decomposition must be shown non-reducible to this
  externality account.
  \url{https://www.aeaweb.org/articles?id=10.1257/mic.20200200}
\item
  {[}partial{]} Dirk Bergemann, Alessandro Bonatti, ``Data, Competition,
  and Digital Platforms'', American Economic Review 114(8):2553-2595
  (2024) (2024) --- Nearest formal precedent for locating the
  observer-to-subject value gap in the platform's market position
  (control of match information and on/off-platform access) --- the
  rent-decomposition leg is anticipated in spirit and must be
  differentiated in mechanism.
  \url{https://www.aeaweb.org/articles?id=10.1257/aer.20230478}
\item
  {[}partial{]} Maryam Farboodi, Laura Veldkamp, ``A Model of the Data
  Economy'', NBER WP 28427 / July 2024 version (2021-2024) --- Data as a
  distinct, durable, dynamically depreciating capital stock with
  Bayes-law (information-theoretic) drift --- covers the `time-drifting
  durable data-capital' pair from the observer side; no privacy metric,
  no inalienability, no rent decomposition.
  \url{https://business.columbia.edu/sites/default/files-efs/citation_file_upload/DataGrowth_july2024.pdf}
\item
  {[}partial{]} Imanol Arrieta-Ibarra, Leonard Goff, Diego
  Jimenez-Hernandez, Jaron Lanier, E. Glen Weyl, ``Should We Treat Data
  as Labor? Moving Beyond `Free'\,'', AEA Papers \& Proceedings 108
  (2018) (and Posner \& Weyl, Radical Markets, 2018) (2018) --- Prices a
  subject-attached, non-transferable-in-stock factor (data-as-labor flow
  with dividends to subjects) --- prior art for compensating subjects
  while the underlying source remains attached to the person,
  compressing the `subject-owned, priced, non-alienated' framing; no
  reconstruction bound, no adversary model, no drift mechanics.
  \url{https://www.aeaweb.org/conference/2018/preliminary/paper/2Y7N88na}
\item
  {[}partial{]} Ignacio N. Cofone, ``Beyond Data Ownership'', 43 Cardozo
  Law Review (2021) (2021) --- Maps data rights onto the
  Calabresi-Melamed property/liability/inalienability grid and argues
  propertization fails on asymmetric information and inferred-data
  leakage --- the structural-inalienability leg is occupied
  conceptually; the claim must show why `inalienable durable capital'
  evades these objections (alongside Malgieri \& Custers, ``Pricing
  privacy -- the right to know the value of your personal data'', CLSR
  2018).
  \url{https://papers.ssrn.com/sol3/papers.cfm?abstract_id=3564480}
\item
  {[}partial{]} Meng Zhang, Ermin Wei, Randall Berry, Jianwei Huang,
  ``Age-Dependent Differential Privacy'', IEEE Trans. Information Theory
  / arXiv:2209.01466 (2022) (2022) --- A formally time-indexed privacy
  guarantee (privacy as an explicit function of data age), alongside the
  continual-observation strand (Dwork et al.~2010; Cao et al.~temporal
  privacy leakage arXiv:1711.11436; temporally discounted DP
  arXiv:1908.03995) --- `first time-indexed privacy quantity' is
  unavailable as a novelty basis; drift direction and absence of pricing
  are the only distinguishers. \url{https://arxiv.org/abs/2209.01466}
\item
  {[}partial{]} Ruwimal Y. Pathiraja, Jerome P. Reiter, ``Setting the
  Privacy Budget in Differential Privacy by Bounding Adversaries' Odds
  of Learning Sensitive Information'', arXiv:2607.04004 (2026) (2026)
  --- Makes the operative privacy budget adversary-relative (calibrated
  to a specified adversary's posterior odds) --- the
  adversary-relativity leg of the budget-setting step is now occupied in
  the DP mainstream; static and unpriced.
  \url{https://arxiv.org/abs/2607.04004}
\item
  {[}partial{]} Joydeep Chandra, ``CHRONOS: Temporally-Aware Multi-Agent
  Coordination for Evolving Data Marketplaces'', arXiv:2605.23887 (2026)
  (2026) --- Joins time-drifting valuation and privacy-budget
  consumption in one marketplace system --- with Bo et
  al.~(arXiv:2512.18296, 2512.21585) shows the DP-markets strand
  actively adding dynamics and strategic pricing; privacy object remains
  a DP budget, no inalienability or rent decomposition.
  \url{https://arxiv.org/abs/2605.23887}
\end{itemize}

\textbf{Absence report (D2).} not\_found\_in\_scope(a single work
jointly pricing an adversary-relative reconstruction bound that drifts
over time via adversary background-information accumulation AND holding
it as inalienable subject-owned durable capital AND decomposing the
observer-to-subject value gap as market-position/venue-exclusion rent).
Searched (8 distinct WebSearch queries beyond the sweep digest, plus
targeted WebFetch of arXiv:2101.04357): privacy depreciation/decay
pricing with adversary knowledge accumulation; data-as-labor/inalienable
data rights pricing (Posner-Weyl lineage); dynamic anonymity pricing
under repeat purchases (Conitzer-Taylor-Wagman lineage);
adversary-dependent privacy pricing with time-varying risk and
reconstruction bounds; information-rent/platform-intermediary
decompositions of the personal-data value gap; inalienable-right and
human-capital pricing analogies for personal data; continual-observation
privacy-loss accumulation with compensation/budget-depletion pricing;
and reconstruction-bound-priced data valuation 2025-2026. Every hit
covers at most one leg or one pair (drift+pricing: Zhang-Zhu 2021, Xu et
al.~2017, CHRONOS 2026; adversary-relativity+reconstruction: Balle et
al.~2022 / ReRo line; drift+durable-capital: Farboodi-Veldkamp;
inalienability+compensation: data-as-labor, Cofone, Malgieri-Custers;
gap decomposition: Bergemann-Bonatti 2024, Acemoglu et al.~2022).
Absence of a joint-quadruple hit is a scope-limited search result, not
proof of novelty.

\subsubsection{CTR-LR-03 --- NARROWED (moderate-high,
75\%)}\label{ctr-lr-03-narrowed-moderate-high-75}

\textbf{Residual restatement (what now survives).} Survives only as the
specific identification: the measured WTA/WTP floor and the IEEE
7012-2025 proposer inversion as TWO EVALUATIONS OF ONE preference-free
subgame-perfect share function of proposer identity, with the subject
share written as an explicit venue-exclusion parameter s = h\_X.
Stripped: any claim to first ultimatum-game treatment of privacy
(Rajtmajer-Squicciarini 2017); the SPE-share-of-proposer mathematics
(Stahl-Rubinstein, known art); and the structural-not-cognitive floor
thesis AS SUCH, which is formally occupied by Gu 2024 (intermediary
bargaining power) and Fallah et al.~2024 (competition benefits buyers
not users) and qualitatively by Zuiderveen Borgesius 2017. Added burden:
an explicit identification argument separating venue exclusion from the
choice-architecture/reference-dependence rival (Lin \& Strulov-Shlain
2023; arXiv:2607.27239; Fradkin), which explains large floors with no
market-structure parameter. The 7012-as-permutation formalization itself
remains unclaimed --- no indexed work links MyTerms/7012 to any
game-theoretic analysis, and no Stackelberg privacy-pricing variant
makes the subject the leader.

\textbf{Basis.} Two refuter/scout scopes ran 9-12 targeted queries
including direct 7012+game-theory and role-reversal Stackelberg searches
with zero hits (a meaningful dog-that-didn't-bark given the standard
published January 2026); but the structural-floor component is
independently formalized (Gu, Fallah) and the empirical floor has a
strong live cognitive rival, so the residue is genuinely smaller than
stated. Strongest surviving residue of the five.

\textbf{Covering art found by the refuter seat.}

\begin{itemize}
\tightlist
\item
  {[}partial{]} Sarah Rajtmajer, Anna Squicciarini et al., ``An
  Ultimatum Game Model for the Evolution of Privacy in Jointly Managed
  Content'', Decision and Game Theory for Security (GameSec), Springer
  LNCS (2017) --- Prior art applying the ultimatum-game formalism inside
  the privacy domain (co-managed content negotiation). It does not touch
  WTA/WTP, venue exclusion, or proposer inversion, but it destroys any
  implicit claim to being the first ultimatum-game treatment of privacy;
  the residue must be stated as the specific WTA/WTP-plus-7012
  identification, not as `ultimatum model of privacy'.
  \url{https://link.springer.com/chapter/10.1007/978-3-319-68711-7_7}
\item
  {[}partial{]} Dale O. Stahl, ``Subgame Perfection in Ultimatum
  Bargaining Trees'' (Stahl--Rubinstein lineage: SPE shares as a
  function of the bargaining tree / who proposes) (c.~2007) --- The
  formal engine of the claim --- that the SPE share is a function of
  proposer identity / move order, and that permuting the proposer
  carries the responder from the floor to an interior share --- is
  textbook Stahl/Rubinstein content, explicitly analyzed across
  bargaining-tree structures. The claim cannot own the
  share-function-of-proposer-identity mathematics; only the
  privacy-domain mapping is left.
  \url{https://laits.utexas.edu/~stahl/experimental/ult_trees.pdf}
\item
  {[}partial{]} Jiadong Gu, ``Data Trade and Consumer Privacy'',
  arXiv:2406.12457 (econ.TH) (2024 (v3 2026)) --- Formal result that the
  data subject's realized privacy/surplus is governed by the
  intermediary's bargaining power, and that reducing that power protects
  the subject --- a structural market-position determinant of the
  subject share, covering the `floor is structural not cognitive' thesis
  in a formal model, though via screening rather than an SPE share
  function of proposer identity. \url{https://arxiv.org/abs/2406.12457}
\item
  {[}partial{]} Alireza Fallah, Michael I. Jordan, Ali Makhdoumi,
  Azarakhsh Malekian, ``On Three-Layer Data Markets'', arXiv:2402.09697
  (2024 (v4 2025)) --- Equilibrium analysis in which the subject's low
  share is a consequence of market structure/position (competition
  benefits buyers, not users) --- independent formal cover for the
  structural (non-preferential) account of the subject's floor, without
  ultimatum unification or a proposer-inversion remedy.
  \url{https://arxiv.org/abs/2402.09697}
\item
  {[}partial{]} Frederik Zuiderveen Borgesius et al., ``Tracking Walls,
  Take-It-Or-Leave-It Choices, the GDPR, and the ePrivacy Regulation'',
  European Data Protection Law Review (re-hosted arXiv:2510.25339)
  (2017) --- Explicit legal-economic treatment of privacy consent as a
  take-it-or-leave-it offer with no negotiation venue --- the
  qualitative structural reading of the floor (consent not freely given
  because the subject cannot counter-offer), predating and paralleling
  the Richards--Hartzog lineage; no SPE model, no share function, no
  h\_X. \url{https://arxiv.org/pdf/2510.25339}
\item
  {[}partial{]} Lijun Bo \& Weiqiang Chang, ``Privacy Data Pricing: A
  Stackelberg Game Approach'', arXiv:2512.18296 (2025) --- Formalizes
  leader/proposer identity in a privacy-pricing game --- closest formal
  object to proposer-identity parameterization; my role-reversal search
  confirmed no variant in this lineage makes the subject the leader, so
  the permutation leg remains open.
  \url{https://arxiv.org/abs/2512.18296}
\item
  {[}partial{]} Tesary Lin \& Avner Strulov-Shlain, ``Choice
  Architecture, Privacy Valuations, and Selection Bias in Consumer
  Data'' (2023 (rev. 2024)) --- Strongest empirical statement of the
  rival cognitive/architectural account of the measured WTA/WTP floor;
  it covers the empirical-floor leg of the claim and forces the
  structural (s = h\_X) reading to carry an explicit identification
  argument distinguishing venue exclusion from elicitation architecture.
  \url{https://arxiv.org/pdf/2308.13496}
\item
  {[}partial{]} IEEE Standards Association, ``IEEE 7012-2025 ---
  Standard for Machine Readable Personal Privacy Terms'' (published
  January 2026) (2026) --- The proposer-inversion instrument exists as a
  ratified artifact self-described as making the individual first party;
  covers the remedy leg as engineering/advocacy fact, with no economic
  model of why the inversion moves the share --- leaving only the
  formalization open.
  \url{https://ieeexplore.ieee.org/document/11360682}
\end{itemize}

\textbf{Absence report (D2).} not\_found\_in\_scope(any published work
casting the empirically measured privacy WTA/WTP floor and the IEEE
7012/MyTerms proposer inversion as two evaluations of ONE
preference-free subgame-perfect-equilibrium share function of proposer
identity, or writing the subject share as a structural venue-exclusion
parameter s = h\_X). Searched 2026-08-17, nine distinct web queries
beyond the sweep digest: (1) ultimatum game + privacy WTA/WTP
bargaining; (2) take-it-or-leave-it privacy consent + subgame perfect +
data market subject share; (3) IEEE 7012/MyTerms + economics/game
theory/bargaining-power/first-party; (4) privacy paradox WTA/WTP gap +
structural bargaining-power (vs preference) explanations; (5) Rubinstein
bargaining + personal data + outside option + proposer role; (6) data
subjects as ultimatum responders / platform as proposer formal models;
(7) ultimatum game analogies to privacy policies / adhesion contracts;
(8) MyTerms/IEEE P7012 + ultimatum/bargaining/equilibrium (zero hits
linking the standard to any game-theoretic analysis); (9) Stackelberg
privacy pricing with user/consumer as leader (role reversal) --- none
found; the entire Stackelberg data-pricing lineage keeps the
platform/collector as leader. Also searched missing-market/Coase
structural accounts of the privacy valuation gap and
consent-as-ultimatum law-review framings. Absence is reported as absence
in this scope, not as proof of novelty.

\subsubsection{CTR-LR-04 --- NARROWED (moderate,
68\%)}\label{ctr-lr-04-narrowed-moderate-68}

\textbf{Residual restatement (what now survives).} The DESTINATION is
occupied: `the redistributive object is position rent, not a wage' is
stated nearly verbatim in Mohades-Savona 2026 (marginal product
attributed to data owners reflects monopoly position, via Marxian rent
theory; early-stage draft --- track to publication before relying on
priority), and is convergently reached by Sadowski 2019,
O'Reilly-Strauss-Mazzucato 2024 (attention rents), and Viljoen 2021, who
runs the structurally isomorphic internal-defeat-then-respecify move.
Acemoglu 2022 supplies a formal near-zero-marginal-price result and
Vincent-Hecht-Sen 2019 ran the diminishing-returns substrate against
data strikes but patched the programme. What survives: the specific
DIRECTED falsifier route --- Bajari et al.~2019 flat returns-to-scale
deployed as an internal empirical defeat of the marginal-product wage
logic, concluding self-defeat rather than patching --- landing on
venue-access/clearing-position rent INSIDE a reconstruction-bound / IEEE
7012 frame. The paper must position this route against three
convergent-but-distinct falsifiers (Mohades-Savona: Marxian rent;
Acemoglu: substitutability; Viljoen: relationality) and against the
competing liability re-specification (Zhang-Zhu 2025).

\textbf{Basis.} Refuter did a full-text read of the strong candidate
plus 9 targeted queries; the conclusion-space is demonstrably crowded
(one strong + four partial candidates), leaving only the argument-route
composition. Confidence is moderate rather than high because
Mohades-Savona is an unpublished draft whose final form could expand to
cover more of the route, and because route-versus-destination
distinctions are the kind reviewers discount.

\textbf{Covering art found by the refuter seat.}

\begin{itemize}
\tightlist
\item
  {[}strong{]} Siavash Mohades \& Maria Savona, ``Data Rent and Surplus
  Value in the Digital Economy'', LUISS LEAP / Maastricht / Sussex
  working paper (early-stage draft, do not cite), January 19, 2026
  (2026) --- Directly occupies the claim's re-specification leg: engages
  Posner-Weyl's data-as-labor by name, formally distinguishes
  differential vs absolute DATA RENT (adapted from Marx's ground rent)
  arising from monopolistic control of access, and states explicitly
  that ``the `marginal product' that neoclassical theory attributes to
  data owners reflects not their productive contribution but their
  monopoly position'' and that ``data rent is not a payment for data's
  `marginal product'\,''. That is the claim's conclusion ---
  redistributive object = position rent, not a wage --- published before
  WP-14. It reaches it via Marxian rent theory rather than the Bajari
  flat-returns falsifier, and its remedy is data sharing/commons rather
  than subject-side venue-access clearing position; no reconstruction
  bound or IEEE 7012.
  \url{https://leap.luiss.it/wp-content/uploads/2026/02/mohades_paper.pdf}
\item
  {[}partial{]} Nicholas Vincent, Brent Hecht, Shilad Sen, ``\,`Data
  Strikes': Evaluating the Effectiveness of a New Form of Collective
  Action Against Technology Companies'', The Web Conference (WWW) 2019
  (2019) --- The closest prior art to RUNNING the falsifier: an
  internal, empirical application of diminishing returns-to-data to the
  data-labour lever itself, simulating data strikes on recommender
  systems and showing small withdrawals have limited effect precisely
  because marginal records add little. Same year and same empirical
  substrate as Bajari 2019. It stops short of the claim's conclusion ---
  it concludes strikes need moderate scale (patching the programme)
  rather than declaring the wage remedy self-defeating and re-specifying
  the object as rent. Companion CDC paper (CSCW 2021,
  dl.acm.org/doi/10.1145/3449177) quantifies the same
  diminishing-returns constraint.
  \url{https://dl.acm.org/doi/10.1145/3308558.3313742}
\item
  {[}partial{]} Salome Viljoen, ``A Relational Theory of Data
  Governance'', Yale Law Journal 131(2) (2021); popular version ``Data
  as Property?'', Phenomenal World (2020) (2021) --- Structurally
  isomorphic falsifier-plus-re-specification: argues the
  propertarian/data-as-labor ``market rate for personal data'' remedy is
  mis-specified because data's value is population-level and not
  reducible to individual contributions, and re-specifies the
  redistributive/governance object away from individual wages. Mechanism
  differs (relationality/sociality, not flat returns-to-scale evidence)
  and the landing differs (democratic collective governance, not
  venue-access rent), but the argumentative move --- internal defeat of
  the wage frame forcing re-specification of the object --- is prior
  art.
  \url{https://yalelawjournal.org/feature/a-relational-theory-of-data-governance}
\item
  {[}partial{]} Daron Acemoglu, Ali Makhdoumi, Azarakhsh Malekian, Asu
  Ozdaglar, ``Too Much Data: Prices and Inefficiencies in Data
  Markets'', American Economic Journal: Microeconomics 14(4): 218-256
  (2022) (2022) --- Formalizes why the price of an individual's marginal
  record is structurally depressed toward zero (data externalities /
  substitutability across correlated users), implying no
  bargaining-toward-marginal-product remedy can deliver meaningful
  individual compensation --- the economic core of the self-defeat,
  derived from substitutability rather than flat returns-to-scale, with
  remedies in market regulation/shutdown rather than venue-rent
  redistribution.
  \url{https://www.aeaweb.org/articles?id=10.1257/mic.20200200}
\item
  {[}partial{]} Tim O'Reilly, Ilan Strauss, Mariana Mazzucato,
  ``Algorithmic attention rents: A theory of digital platform market
  power'', Data \& Policy 6 (2024) (2024) --- Prior art for reading
  platform data value as RENT extracted from control of the
  allocation/matching venue (attention allocation position) rather than
  from productive contribution --- the same clearing-position-rent
  object the claim re-specifies toward, without any engagement with the
  data-as-labour wage remedy or the flat-returns falsifier.
  \url{https://www.cambridge.org/core/journals/data-and-policy/article/algorithmic-attention-rents-a-theory-of-digital-platform-market-power/D85FE41F6CF99FC57DDFB2B2B63491C5}
\end{itemize}

\textbf{Absence report (D2).} not\_found\_in\_scope(9 web searches +
full-text read of Mohades-Savona 2026 across: {[}"data as labor"
critique diminishing returns/marginal value wage self-defeating{]},
{[}Bajari returns-to-scale + data-as-labor/data-dividend critique{]},
{[}data dividend criticism marginal value near zero{]}, {[}Acemoglu Too
Much Data substitutability + data unions{]}, {[}Viljoen relational
theory vs propertarian compensation{]}, {[}Sadowski data-as-capital
rent{]}, {[}Savona data rents redistribution vs wages{]},
{[}Vincent-Hecht data strikes diminishing returns{]}, {[}Posner-Weyl
marginal product critique{]}, {[}self-defeating data-dividend +
rent-instead phrasing{]}): no work was found that (a) deploys Bajari et
al.~2019 flat returns-to-scale evidence AS a directed falsifier
concluding the data-as-labour wage remedy self-defeats, AND (b)
re-specifies the redistributive object as venue-access/clearing-position
rent in the same chain, AND (c) does either inside a
reconstruction-bound or IEEE 7012 frame. Absence of this composition is
reported as search-scope absence, not as proof of novelty.

\subsubsection{CTR-LR-05 --- NARROWED (moderate,
72\%)}\label{ctr-lr-05-narrowed-moderate-72}

\textbf{Residual restatement (what now survives).} Stripped: `first to
make auxiliary/background-information accumulation the shelf-life clock'
(Altman et al.~2018 do it in governance form; Rocher et al.~2019
quantitatively; Acemoglu 2022 economically via the correlation
externality); `a durability/depreciation term for data-as-capital'
(Veldkamp 2023; Farboodi-Veldkamp; NBER w35266 endogenous depreciation;
Coyle-Manley 2024; Valavi 2022); and `Mosca-style shelf-life reasoning
in economics venues' (Mascelli-Rodden FEDS 2025-093; arXiv:2603.01091).
Survives as the two-part composition: (a) writing a Mosca-FORM
inequality whose clock is an information-theoretic reconstruction bound
(Dinur-Nissim / QIF / Bayes-capacity sense, e.g.~arXiv:2406.13569) in
place of cryptographic break-time, AND (b) carrying that substituted
clock into a SUBJECT-OWNED capital-asset valuation --- with the clock
explicitly distinguished from Veldkamp-style relevance-decay on both
axes (adversary-side accumulation vs statistical obsolescence;
subject-owned vs firm-owned asset).

\textbf{Basis.} Sixteen distinct queries across two scopes found the
ingredient (QIF reconstruction bounds), the template (Mosca as portable
practitioner framework), and the destination class (data-capital
depreciation) separately but never the substitution or its valuation
carriage; however Altman 2018 pre-empts the conceptual half and the
practitioner generalization of Mosca weakens the inventiveness of
re-templating, so only the formal composition survives. One flagged
unresolved item (ResearchGate-only time-dependent HNDL framework) needs
manual follow-up before submission.

\textbf{Covering art found by the refuter seat.}

\begin{itemize}
\tightlist
\item
  {[}strong{]} Laura Veldkamp, ``Valuing Data as an Asset'', Review of
  Finance 27(5); Maryam Farboodi, Laura Veldkamp, ``A Model of the Data
  Economy'' (NBER WP 28427; July 2024 version) (2021-2024) --- Covers
  the broader claim shape as the digest already flags: data as
  long-lived, depreciating capital with a formal,
  information-dynamics-driven depreciation rate (Bayes/Kalman precision
  decay). Firm-owned and obsolescence-clocked, not adversary-clocked or
  subject-owned --- but it is the strongest covering art for `a
  durability term for data-as-capital' as such.
  \url{https://business.columbia.edu/sites/default/files-efs/citation_file_upload/DataGrowth_july2024.pdf}
\item
  {[}partial{]} Micah Altman, Alexandra Wood, David O'Brien, Salil
  Vadhan, Urs Gasser, ``Practical approaches to big data privacy over
  time'', International Data Privacy Law 8(1) (2018) --- Anticipates the
  conceptual half of the substitution: protection afforded by
  SDL/de-identification has a finite horizon set by accumulating
  auxiliary (background) information, and governance should be designed
  over the data lifetime. This is `background-information accumulation
  as the durability clock' in governance form --- it lacks the
  Mosca-form inequality and any capital-asset valuation, but it strips
  the claim of `first to make auxiliary-info accumulation the shelf-life
  clock'. \url{https://academic.oup.com/idpl/article/8/1/29/4930711}
\item
  {[}partial{]} Luc Rocher, Julien M. Hendrickx, Yves-Alexandre de
  Montjoye, ``Estimating the success of re-identifications in incomplete
  datasets using generative models'', Nature Communications 10:3069
  (2019) --- Formal, quantitative model of re-identification success
  showing anonymization is not future-proof as datasets accumulate ---
  an operationalized, adversary-relative reconstruction-risk quantity
  that grows with available background data. Attack-side and
  governance-facing; no inequality, no valuation, but it is the standard
  citation a reviewer would raise against `a finite, adversary-relative,
  drifting reconstruction bound' being new.
  \url{https://www.nature.com/articles/s41467-019-10933-3}
\item
  {[}partial{]} Chao Li, Daniel Yang Li, Gerome Miklau, Dan Suciu, ``A
  Theory of Pricing Private Data'', ICDT 2013 / arXiv:1208.5258 (with
  Ghosh-Roth 2011 auction lineage) (2012) --- Covers the subject-owned
  valuation half in static form: data subjects are compensated in
  proportion to an information-theoretic privacy-loss quantity
  (differential-privacy leakage) --- i.e., an information-theoretic
  bound already prices subject-owned data. No time clock and no
  Mosca-form inequality, but it blocks any reading of the claim as
  `first to price subject-owned data by an information-theoretic bound'.
  \url{https://arxiv.org/pdf/1208.5258}
\item
  {[}partial{]} Justin Hsu, Marco Gaboardi, Andreas Haeberlen, Sanjeev
  Khanna, Arjun Narayan, Benjamin C. Pierce, Aaron Roth, ``Differential
  Privacy: An Economic Method for Choosing Epsilon'', IEEE CSF 2014 /
  arXiv:1402.3329 (2014) --- Bridges an information-theoretic privacy
  bound directly to economic compensation owed to data subjects for the
  risk they bear --- the bound-to-subject-economics bridge, static in
  time. Reinforces that only the temporal/inequality composition, not
  the economics-of-the-bound, is residually open.
  \url{https://arxiv.org/pdf/1402.3329}
\item
  {[}partial{]} Lin William Cong, Danxia Xie, Longtian Zhang,
  ``Knowledge Accumulation, Privacy, and Growth in a Data Economy'',
  Management Science 68(9) / arXiv:2109.10028 (2022) ---
  Consumer-generated (subject-side) data enters a
  capital/knowledge-accumulation growth model with an explicit
  privacy-infringement cost that evolves dynamically ---
  subject-provided data as a dynamic factor with a time-path privacy
  cost. The privacy cost is a preference/externality term, not a
  reconstruction bound, and there is no shelf-life inequality; but it
  occupies the `subject-side data capital with dynamic privacy cost'
  ground. \url{https://arxiv.org/abs/2109.10028}
\item
  {[}partial{]} Daron Acemoglu, Ali Makhdoumi, Azarakhsh Malekian,
  Asuman Ozdaglar, ``Too Much Data: Prices and Inefficiencies in Data
  Markets'', AEJ: Microeconomics 14(4) (2022) --- The data externality
  --- the informativeness of others' accumulated data erodes an
  individual's privacy and depresses the price of their data --- is
  economically the same mechanism as `background-information
  accumulation degrades the subject's asset'. Never cast as a shelf-life
  inequality or durability term, but it pre-occupies the economic
  content of adversary-side accumulation lowering subject-side data
  value. \url{https://www.aeaweb.org/articles?id=10.1257/mic.20200200}
\item
  {[}partial{]} Meng Zhang, Ermin Wei, Randall Berry, Jianwei Huang,
  ``Age-Dependent Differential Privacy'', IEEE Trans. Information Theory
  / arXiv:2209.01466 (2022) --- A formally time-indexed
  information-theoretic privacy quantity (privacy risk as a function of
  data age) --- the same mathematical object class as the substituted
  clock, applied to utility trade-offs rather than asset valuation and
  without a Mosca-form inequality.
  \url{https://arxiv.org/abs/2209.01466}
\item
  {[}partial{]} Jillian Mascelli, Megan Rodden, ``Harvest Now Decrypt
  Later: Examining Post-Quantum Cryptography and the Data Privacy Risks
  for Distributed Ledger Networks'', Federal Reserve FEDS WP 2025-093
  (2025) --- Closest published move of Mosca-style shelf-life reasoning
  into data-privacy economics (retrospective confidentiality horizon on
  recorded personal/financial data at an economics institution). Clock
  remains cryptographic break-time; no reconstruction-bound
  substitution, no subject-owned asset valuation.
  \url{https://www.federalreserve.gov/econres/feds/files/2025093pap.pdf}
\end{itemize}

\textbf{Absence report (D2).} not\_found\_in\_scope(8 WebSearch sweeps,
2026-08-17: (1) `shelf life' + reconstruction attack + auxiliary
information + time-horizon inequality; (2) Mosca-inequality analog for
de-identification/anonymization lifetime under auxiliary-information
accumulation; (3) privacy half-life / re-identification value-decay
economics; (4) Dinur-Nissim reconstruction bound + economic
valuation/pricing of personal data over time; (5) `harvest now,
re-identify/deanonymize later' for anonymized personal data; (6)
`privacy depreciation' as an economic asset model driven by adversary
information accumulation; (7) Mosca theorem cross-cited with
differential privacy/reconstruction/re-identification; (8) subject-owned
inalienable data-as-capital valuation with adversary-indexed durability.
No source was found that (a) writes a Mosca-FORM inequality
(protection-lifetime + migration/exposure time vs clock) whose clock is
an information-theoretic reconstruction bound (Dinur-Nissim/QIF) rather
than cryptographic break-time, NOR (b) carries such a substituted clock
into a valuation formula for a subject-owned capital asset. Absence of
found art is reported as search outcome only, not as proof of novelty.)

\newpage

\section{Bibliography additions proposed by run
03}\label{bibliography-additions-proposed-by-run-03}

74 new resolvable citations surfaced this run and judged worth adding to
the WP-14 bibliography. Provenance-gated per run-02 practice: A4
(citation-verifier) re-verifies every entry against a primary record
before any tier-A use; \texttt{provenance\_confidence} below is the
runtime's own label, not A4's.

\begin{enumerate}
\def\labelenumi{\arabic{enumi}.}
\tightlist
\item
  Vincent Conitzer, Curtis R. Taylor, Liad Wagman, ``Hide and Seek:
  Costly Consumer Privacy in a Market with Repeat Purchases'', Marketing
  Science 31(2):277-292 (2012).
  \url{https://papers.ssrn.com/sol3/papers.cfm?abstract_id=2160895} ·
  provenance\_confidence: high
\item
  Justin Hsu, Marco Gaboardi, Andreas Haeberlen, Sanjeev Khanna, Arjun
  Narayan, Benjamin C. Pierce, Aaron Roth, ``Differential Privacy: An
  Economic Method for Choosing Epsilon'', IEEE CSF 2014, arXiv:1402.3329
  (2014). \url{https://arxiv.org/abs/1402.3329} ·
  provenance\_confidence: high
\item
  Frederik Zuiderveen Borgesius et al., ``Tracking Walls,
  Take-It-Or-Leave-It Choices, the GDPR, and the ePrivacy Regulation'',
  European Data Protection Law Review (2017).
  \url{https://arxiv.org/pdf/2510.25339} · provenance\_confidence:
  medium
\item
  Lei Xu et al., ``Dynamic Privacy Pricing: A Multi-Armed Bandit
  Approach With Time-Variant Rewards'', IEEE Trans. Information
  Forensics and Security (2017).
  \url{https://www.researchgate.net/publication/316236814_Dynamic_Privacy_Pricing_A_Multi-Armed_Bandit_Approach_With_Time-Variant_Rewards}
  · provenance\_confidence: medium
\item
  Sarah Rajtmajer, Anna Squicciarini et al., ``An Ultimatum Game Model
  for the Evolution of Privacy in Jointly Managed Content'', GameSec
  2017, Springer LNCS (2017).
  \url{https://link.springer.com/chapter/10.1007/978-3-319-68711-7_7} ·
  provenance\_confidence: high
\item
  Kobbi Nissim, Aaron Bembenek, Alexandra Wood, Mark Bun, Marco
  Gaboardi, Urs Gasser, David R. O'Brien, Salil Vadhan, Thomas Steinke,
  ``Bridging the Gap between Computer Science and Legal Approaches to
  Privacy'', Harvard Journal of Law \& Technology 31(2) (2018).
  \url{https://www.semanticscholar.org/paper/Bridging-the-Gap-between-Computer-Science-and-Legal-Nissim-Bembenek/22067befe017dc5760be667bf0ba9cf48b2ab650}
  · provenance\_confidence: high
\item
  Micah Altman, Alexandra Wood, David O'Brien, Salil Vadhan, Urs Gasser,
  ``Practical approaches to big data privacy over time'', International
  Data Privacy Law 8(1) (2018).
  \url{https://academic.oup.com/idpl/article/8/1/29/4930711} ·
  provenance\_confidence: high
\item
  Cynthia Dwork, Nitin Kohli, Deirdre Mulligan, ``Differential Privacy
  in Practice: Expose your Epsilons!'', Journal of Privacy and
  Confidentiality 9(2) (2019).
  \url{https://journalprivacyconfidentiality.org/index.php/jpc/article/view/689}
  · provenance\_confidence: high
\item
  Jathan Sadowski, ``When data is capital: Datafication, accumulation,
  and extraction'', Big Data \& Society 6(1) (2019).
  \url{https://journals.sagepub.com/doi/10.1177/2053951718820549} ·
  provenance\_confidence: high
\item
  John M. Abowd, Ian M. Schmutte, ``An Economic Analysis of Privacy
  Protection and Statistical Accuracy as Social Choices'', American
  Economic Review 109(1), arXiv:1808.06303 (2019).
  \url{https://arxiv.org/abs/1808.06303} · provenance\_confidence: high
\item
  Luc Rocher, Julien M. Hendrickx, Yves-Alexandre de Montjoye,
  ``Estimating the success of re-identifications in incomplete datasets
  using generative models'', Nature Communications 10:3069 (2019).
  \url{https://www.nature.com/articles/s41467-019-10933-3} ·
  provenance\_confidence: high
\item
  Nicholas Vincent, Brent Hecht, Shilad Sen, ``\,`Data Strikes':
  Evaluating the Effectiveness of a New Form of Collective Action
  Against Technology Companies'', WWW 2019 (2019).
  \url{https://dl.acm.org/doi/10.1145/3308558.3313742} ·
  provenance\_confidence: high
\item
  Bennett Cyphers, Adam Schwartz, ``Why Getting Paid for Your Data Is a
  Bad Deal'', Electronic Frontier Foundation (2020).
  \url{https://www.eff.org/deeplinks/2020/10/why-getting-paid-your-data-bad-deal}
  · provenance\_confidence: high
\item
  U.S. Census Bureau (Abowd et al.), ``The modernization of statistical
  disclosure limitation at the U.S. Census Bureau'', working paper
  (2020).
  \url{https://www.census.gov/content/dam/Census/library/working-papers/2020/adrm/The\%20modernization\%20of\%20statistical\%20disclosure\%20limitation\%20at\%20the\%20U.S.\%20Census\%20Bureau.pdf}
  · provenance\_confidence: high
\item
  Ignacio N. Cofone, ``Beyond Data Ownership'', 43 Cardozo Law Review
  (2021).
  \url{https://papers.ssrn.com/sol3/papers.cfm?abstract_id=3564480} ·
  provenance\_confidence: high
\item
  Nicholas Vincent, Hanlin Li et al., ``Data Leverage: A Framework for
  Empowering the Public in its Relationship with Technology Companies'',
  FAccT 2021 (2021).
  \url{https://nickmvincent.com/static/data_leverage.pdf} ·
  provenance\_confidence: high
\item
  Salome Viljoen, ``A Relational Theory of Data Governance'', Yale Law
  Journal 131(2) (2021).
  \url{https://yalelawjournal.org/feature/a-relational-theory-of-data-governance}
  · provenance\_confidence: high
\item
  Tao Zhang, Quanyan Zhu, ``On the Differential Private Data Market:
  Endogenous Evolution, Dynamic Pricing, and Incentive Compatibility'',
  arXiv:2101.04357 (2021). \url{https://arxiv.org/abs/2101.04357} ·
  provenance\_confidence: high
\item
  Maryam Farboodi, Laura Veldkamp, ``A Model of the Data Economy'', NBER
  Working Paper 28427 (July 2024 version) (2021-2024).
  \url{https://business.columbia.edu/sites/default/files-efs/citation_file_upload/DataGrowth_july2024.pdf}
  · provenance\_confidence: high
\item
  ``A General Framework for Auditing Differentially Private Machine
  Learning'', NeurIPS 2022, arXiv:2210.08643 (2022).
  \url{https://arxiv.org/pdf/2210.08643} · provenance\_confidence:
  medium
\item
  Borja Balle, Giovanni Cherubin, Jamie Hayes, ``Reconstructing Training
  Data with Informed Adversaries'', IEEE S\&P 2022, arXiv:2201.04845
  (2022). \url{https://arxiv.org/abs/2201.04845} ·
  provenance\_confidence: high
\item
  Daron Acemoglu, Ali Makhdoumi, Azarakhsh Malekian, Asuman Ozdaglar,
  ``Too Much Data: Prices and Inefficiencies in Data Markets'', AEJ:
  Microeconomics 14(4):218-256 (2022).
  \url{https://www.aeaweb.org/articles?id=10.1257/mic.20200200} ·
  provenance\_confidence: high
\item
  Ehsan Valavi, Joel Hestness, Newsha Ardalani, Marco Iansiti, ``Time
  and the Value of Data'', HBS WP, arXiv:2203.09118 (2022).
  \url{https://arxiv.org/pdf/2203.09118} · provenance\_confidence:
  medium
\item
  Lin William Cong, Danxia Xie, Longtian Zhang, ``Knowledge
  Accumulation, Privacy, and Growth in a Data Economy'', Management
  Science 68(9), arXiv:2109.10028 (2022).
  \url{https://arxiv.org/abs/2109.10028} · provenance\_confidence: high
\item
  Meng Zhang, Ermin Wei, Randall Berry, Jianwei Huang, ``Age-Dependent
  Differential Privacy'', IEEE Trans. Information Theory,
  arXiv:2209.01466 (2022). \url{https://arxiv.org/abs/2209.01466} ·
  provenance\_confidence: high
\item
  Georgios Kaissis, Jamie Hayes, Alexander Ziller, Daniel Rueckert,
  ``Bounding data reconstruction attacks with the hypothesis testing
  interpretation of differential privacy'', arXiv:2307.03928 (2023).
  \url{https://arxiv.org/abs/2307.03928} · provenance\_confidence: high
\item
  Laura Veldkamp, ``Valuing Data as an Asset'', Review of Finance 27(5)
  (2023).
  \url{https://business.columbia.edu/sites/default/files-efs/citation_file_upload/Valuing\%20Data\%20as\%20an\%20Asset.pdf}
  · provenance\_confidence: high
\item
  Michael Khavkin, Eran Toch, ``Valuation of Differential Privacy Budget
  in Data Trade: A Conjoint Analysis'', Privacy Symposium 2023, Springer
  LNCS (2023).
  \url{https://link.springer.com/chapter/10.1007/978-3-031-31971-6_7} ·
  provenance\_confidence: high
\item
  Tesary Lin, Avner Strulov-Shlain, ``Choice Architecture, Privacy
  Valuations, and Selection Bias in Consumer Data'', arXiv:2308.13496
  (2023). \url{https://arxiv.org/pdf/2308.13496} ·
  provenance\_confidence: high
\item
  ``A Survey on Data Markets'', arXiv:2411.07267 (2024).
  \url{https://arxiv.org/html/2411.07267} · provenance\_confidence:
  medium
\item
  Alessandro Acquisti, ``The Economics of Privacy at a Crossroads'', in
  Goldfarb \& Tucker (eds.), The Economics of Privacy, NBER/Univ. of
  Chicago Press (2024).
  \url{https://papers.ssrn.com/sol3/papers.cfm?abstract_id=4898619} ·
  provenance\_confidence: high
\item
  Alireza Fallah, Ali Makhdoumi, Azarakhsh Malekian, Asuman Ozdaglar,
  ``Optimal and Differentially Private Data Acquisition: Central and
  Local Mechanisms'', Operations Research 72(3):1105-1123,
  arXiv:2201.03968 (2024). \url{https://arxiv.org/abs/2201.03968} ·
  provenance\_confidence: high
\item
  Alireza Fallah, Michael I. Jordan, Ali Makhdoumi, Azarakhsh Malekian,
  ``On Three-Layer Data Markets'', arXiv:2402.09697 (2024).
  \url{https://arxiv.org/abs/2402.09697} · provenance\_confidence: high
\item
  Andrey Fradkin et al., ``Designing Consent: Choice Architecture and
  Consumer Welfare in Data Sharing'', working paper (2024).
  \url{https://andreyfradkin.com/assets/Privacy_Preferences_and_Dark_Patterns.pdf}
  · provenance\_confidence: medium
\item
  Anneliese Riess, Kristian Schwethelm, Johannes Kaiser, Tamara T.
  Mueller, Julia A. Schnabel, Daniel Rueckert, Alexander Ziller, ``From
  Mean to Extreme: Formal Differential Privacy Bounds on the Success of
  Real-World Data Reconstruction Attacks'', arXiv:2402.12861 (2024).
  \url{https://arxiv.org/abs/2402.12861} · provenance\_confidence: high
\item
  Diane Coyle, Annabel Manley, ``What is the value of data? A review of
  empirical methods'', Journal of Economic Surveys 39(2) (2024).
  \url{https://onlinelibrary.wiley.com/doi/full/10.1111/joes.12585} ·
  provenance\_confidence: high
\item
  Diptangshu Sen, Jingyan Wang, Juba Ziani, ``Equilibria of Data
  Marketplaces with Privacy-Aware Sellers under Endogenous Privacy
  Costs'', arXiv:2402.08826 (2024).
  \url{https://arxiv.org/abs/2402.08826} · provenance\_confidence: high
\item
  Dirk Bergemann, Alessandro Bonatti, ``Data, Competition, and Digital
  Platforms'', American Economic Review 114(8):2553-2595 (2024).
  \url{https://www.aeaweb.org/articles?id=10.1257/aer.20230478} ·
  provenance\_confidence: high
\item
  Jiadong Gu, ``Data Trade and Consumer Privacy'', arXiv:2406.12457
  (econ.TH) (2024). \url{https://arxiv.org/abs/2406.12457} ·
  provenance\_confidence: high
\item
  Lina Dencik, Jessica Brand, Sarah Murphy, ``What do data rights do for
  workers? A critical analysis of trade union engagement with the
  datafied workplace'', Transfer: European Review of Labour and Research
  (2024).
  \url{https://journals.sagepub.com/doi/10.1177/10242589241267006} ·
  provenance\_confidence: high
\item
  Roxana Mihet, Orlando Gomes, Kumar Rishabh, ``Data Risk, Firm Growth
  and Innovation'', WEIS 2024 (2024).
  \url{https://weis.utdallas.edu/program} · provenance\_confidence:
  medium
\item
  Sara Saeidian et al., ``Information Density Bounds for Privacy'',
  arXiv:2407.01167 (2024). \url{https://arxiv.org/pdf/2407.01167} ·
  provenance\_confidence: medium
\item
  Saurab Chhachhi, Fei Teng, ``Wasserstein Markets for
  Differentially-Private Data'', arXiv:2412.02609 (2024).
  \url{https://arxiv.org/abs/2412.02609} · provenance\_confidence: high
\item
  Sayan Biswas, Mark Dras, Pedro Faustini, Natasha Fernandes, Annabelle
  McIver, Catuscia Palamidessi, Parastoo Sadeghi, ``Bayes' capacity as a
  measure for reconstruction attacks in federated learning'',
  arXiv:2406.13569 (2024). \url{https://arxiv.org/abs/2406.13569} ·
  provenance\_confidence: high
\item
  Shuaiqi Wang et al., ``Statistic Maximal Leakage'', arXiv:2411.18531
  (2024). \url{https://arxiv.org/pdf/2411.18531} ·
  provenance\_confidence: medium
\item
  Tim O'Reilly, Ilan Strauss, Mariana Mazzucato, ``Algorithmic attention
  rents: A theory of digital platform market power'', Data \& Policy 6
  (2024).
  \url{https://www.cambridge.org/core/journals/data-and-policy/article/algorithmic-attention-rents-a-theory-of-digital-platform-market-power/D85FE41F6CF99FC57DDFB2B2B63491C5}
  · provenance\_confidence: high
\item
  ``A Unified Framework for Adversary-Aware Differential Privacy
  Bounds'', arXiv:2507.08158 (contents unverified --- fetch before
  citing) (2025). \url{https://arxiv.org/pdf/2507.08158} ·
  provenance\_confidence: medium
\item
  ``Data Cooperatives: Democratic Models for Ethical Data Stewardship'',
  arXiv:2504.10058 (2025). \url{https://arxiv.org/abs/2504.10058} ·
  provenance\_confidence: medium
\item
  ``Learning to Attack: Uncovering Privacy Risks in Sequential Data
  Releases'', arXiv:2510.24807 (2025).
  \url{https://arxiv.org/html/2510.24807v1} · provenance\_confidence:
  high
\item
  ``Multi-Party Data Pricing for Complex Data Trading Markets: A
  Rubinstein Bargaining Approach'', arXiv:2502.16363 (2025).
  \url{https://arxiv.org/html/2502.16363} · provenance\_confidence:
  medium
\item
  ``The Economics of AI Training Data: A Research Agenda'',
  arXiv:2510.24990 (2025). \url{https://arxiv.org/pdf/2510.24990} ·
  provenance\_confidence: medium
\item
  Doc Searls / Customer Commons, ``MyTerms'' (IEEE 7012 first-party
  proffering hub) (2025). \url{https://doc.searls.com/myterms/} ·
  provenance\_confidence: high
\item
  Haoyi Zhang, Tianyi Zhu, ``Neither Consent nor Property: A Policy Lab
  for Data Law'', arXiv:2510.26727 (2025).
  \url{https://arxiv.org/abs/2510.26727} · provenance\_confidence: high
\item
  Jillian Mascelli, Megan Rodden, ``\,`Harvest Now Decrypt Later':
  Examining Post-Quantum Cryptography and the Data Privacy Risks for
  Distributed Ledger Networks'', Federal Reserve FEDS WP 2025-093
  (2025).
  \url{https://www.federalreserve.gov/econres/feds/files/2025093pap.pdf}
  · provenance\_confidence: medium
\item
  Lijun Bo et al., ``Dynamic Data Pricing: A Mean Field Stackelberg Game
  Approach'', arXiv:2512.21585 (2025).
  \url{https://arxiv.org/abs/2512.21585} · provenance\_confidence: high
\item
  Lijun Bo, Weiqiang Chang, ``Privacy Data Pricing: A Stackelberg Game
  Approach'', arXiv:2512.18296 (2025).
  \url{https://arxiv.org/abs/2512.18296} · provenance\_confidence: high
\item
  Luigi D. C. Soares, Mario S. Alvim, Natasha Fernandes, ``A new measure
  for dynamic leakage based on quantitative information flow'',
  arXiv:2510.20922 (2025). \url{https://arxiv.org/abs/2510.20922} ·
  provenance\_confidence: high
\item
  NIST, ``Guidelines for Evaluating Differential Privacy Guarantees'',
  NIST SP 800-226 (final) (2025).
  \url{https://nvlpubs.nist.gov/nistpubs/SpecialPublications/NIST.SP.800-226.pdf}
  · provenance\_confidence: high
\item
  Neil M. Richards, Woodrow Hartzog, Jordan Francis, ``Privacy's
  Autonomy Thicket: Disentangling Choice, Consent, and Control''
  (forthcoming) (2025).
  \url{https://www.private-law-theory.org/2025/12/04/richards-hartzog-and-francis-privacys-autonomy-thicket-disentangling-choice-consent-and-control/}
  · provenance\_confidence: medium
\item
  Nicholas Vincent, Matthew Prewitt, Hanlin Li, ``Collective Bargaining
  in the Information Economy Can Address AI-Driven Power
  Concentration'', arXiv:2506.10272 (2025).
  \url{https://arxiv.org/abs/2506.10272} · provenance\_confidence: high
\item
  Patrick Mesana, Gilles Caporossi, Sebastien Gambs, ``Measuring the
  Hidden Cost of Data Valuation through Collective Disclosure'',
  arXiv:2510.08869 (2025). \url{https://arxiv.org/abs/2510.08869} ·
  provenance\_confidence: high
\item
  Sayan Biswas, Mark Dras, Pedro Faustini, Natasha Fernandes, Annabelle
  McIver, Catuscia Palamidessi, Parastoo Sadeghi, ``Comparing privacy
  notions for protection against reconstruction attacks in machine
  learning'', arXiv:2502.04045 (2025).
  \url{https://arxiv.org/abs/2502.04045} · provenance\_confidence: high
\item
  Woodrow Hartzog, Daniel J. Solove, ``Privacy as Contract?'', Harvard
  Journal of Law \& Technology 39 (forthcoming), SSRN 6081773 (2025).
  \url{https://papers.ssrn.com/sol3/papers.cfm?abstract_id=6081773} ·
  provenance\_confidence: high
\item
  ``Reference Dependence and the Structure of the WTA/WTP Gap'',
  arXiv:2607.27239 (2026). \url{https://arxiv.org/html/2607.27239} ·
  provenance\_confidence: medium
\item
  ``SoK: Reconstruction Attacks on Synthetic Tabular Data (Insights from
  Winning the NIST CRC)'', arXiv:2606.08372 (2026).
  \url{https://arxiv.org/html/2606.08372} · provenance\_confidence:
  medium
\item
  Ankit Bhutani, Guillermo Ordonez, Laura Veldkamp, ``The Missing Value
  of Data'', NBER WP 35266 / AEA 2026 (2026).
  \url{https://www.nber.org/papers/w35266} · provenance\_confidence:
  medium
\item
  IEEE Standards Association, ``IEEE 7012-2025 --- Standard for Machine
  Readable Personal Privacy Terms'' (published January 2026) (2026).
  \url{https://ieeexplore.ieee.org/document/11360682} ·
  provenance\_confidence: high
\item
  Javier Blanco-Romero, Florina Almenares Mendoza, Carlos Garcia Rubio,
  Celeste Campo, Daniel Diaz Sanchez, ``On the Practical Feasibility of
  Harvest-Now, Decrypt-Later Attacks'', arXiv:2603.01091 (2026).
  \url{https://arxiv.org/abs/2603.01091} · provenance\_confidence: high
\item
  Johannes Kaiser, Alexander Ziller, Eleni Triantafillou, Daniel
  Rueckert, Georgios Kaissis, ``Your Privacy Depends on Others:
  Collusion Vulnerabilities in Individual Differential Privacy'', IEEE
  SaTML, arXiv:2601.12922 (2026). \url{https://arxiv.org/abs/2601.12922}
  · provenance\_confidence: high
\item
  Joydeep Chandra, ``CHRONOS: Temporally-Aware Multi-Agent Coordination
  for Evolving Data Marketplaces'', arXiv:2605.23887 (2026).
  \url{https://arxiv.org/abs/2605.23887} · provenance\_confidence: high
\item
  Mohan et al., ``Enterprise Data Valuation --- A Targeted Literature
  Review'', Journal of Economic Surveys (2026).
  \url{https://onlinelibrary.wiley.com/doi/10.1111/joes.12705} ·
  provenance\_confidence: high
\item
  Nielsen, Lev-Aretz et al., ``Privacy Agents in the Field'', WEIS 2026
  accepted paper (no public PDF yet --- pull before WEIS 2027
  submission) (2026).
  \url{https://weis2026.econinfosec.org/accepted-papers/} ·
  provenance\_confidence: low
\item
  Ruwimal Y. Pathiraja, Jerome P. Reiter, ``Setting the Privacy Budget
  in Differential Privacy by Bounding Adversaries' Odds of Learning
  Sensitive Information'', arXiv:2607.04004 (2026).
  \url{https://arxiv.org/abs/2607.04004} · provenance\_confidence: high
\item
  Siavash Mohades, Maria Savona, ``Data Rent and Surplus Value in the
  Digital Economy'', LUISS LEAP working paper (EARLY-STAGE DRAFT marked
  do-not-cite --- track to publication; key LR-04 prior-art threat)
  (2026).
  \url{https://leap.luiss.it/wp-content/uploads/2026/02/mohades_paper.pdf}
  · provenance\_confidence: low
\end{enumerate}

\end{document}
